\documentclass[11pt]{article}

\newcommand{\cnum}{Math 245B}
\newcommand{\ced}{Winter 2026}
\newcommand{\ctitle}[4]{\title{\vspace{-0.5in}\cnum, \ced\\Week #1: #2}\author{\vspace{-0.35in}\\#3, #4}}
\usepackage{enumitem}
\usepackage[usenames,dvipsnames,svgnames,table,hyperref]{xcolor}
\usepackage{amsmath, amsfonts}
\usepackage{graphicx} % Required for inserting images
\usepackage{float}
\usepackage{listings}
\usepackage{xcolor}
\usepackage{hyperref}
\usepackage{comment}

\renewcommand*{\theenumi}{\alph{enumi}}
\renewcommand*\labelenumi{(\theenumi)}
\renewcommand*{\theenumii}{\roman{enumii}}
\renewcommand*\labelenumii{\theenumii.}

\definecolor{codegreen}{rgb}{0,0.6,0}
\definecolor{codegray}{rgb}{0.5,0.5,0.5}
\definecolor{codepurple}{rgb}{0.58,0,0.82}
\definecolor{backcolour}{rgb}{0.95,0.95,0.92}
\definecolor{header}{HTML}{205459}
\definecolor{solution}{HTML}{204d52}
\DeclareMathOperator{\spt}{spt}
\DeclareMathOperator{\iter}{int}

\newcommand{\solution}[1]{{{\textcolor{header}{Solution:} \textcolor{solution}{#1}}}}
\setlist[enumerate]{label=(\alph*)}

\lstdefinestyle{mystyle}{
    backgroundcolor=\color{backcolour},
    commentstyle=\color{codegreen},
    keywordstyle=\color{magenta},
    numberstyle=\tiny\color{codegray},
    stringstyle=\color{codepurple},
    basicstyle=\ttfamily\footnotesize,
    breakatwhitespace=false,
    breaklines=true,
    captionpos=b,
    keepspaces=true,
    numbers=left,
    numbersep=5pt,
    showspaces=false,
    showstringspaces=false,
    showtabs=false,
    tabsize=2
}


\begin{document}
\ctitle{\#2}{}{Asher Christian}{006-150-286}
\date{}
\maketitle

\emph{
    Let $E$ be a a Banach space and denote by $\tau_{E'}$ the strong topology on  $ E'$.
}
Reminder: 
\begin{enumerate}
    \item If $ \mathcal{X}$ is a vector space and $\phi_0, \dots, \phi_n : X \rightarrow \mathbb{R}$ are linear
        such that 
         \[
             \bigcap_{i=1}^{n} Ker(\phi_i) \subset Ker(\phi_0) \implies \phi_0 = \sum_{i=1}^{n}c_i \phi_i
        \] 
    \item We define $J_e : E \rightarrow E''$ by $J_E(x)(f) = f(x)$ if $x \in E$ and  $f \in E'$.
    \item $\sigma(E', E)$ is the topology generated by $\{J_E(x) : x \in E\}$
\end{enumerate}
Remark: If $x \in E$ then  $J_E(x)$ s continuous for the  $\sigma(E',E)$ topology.
Proposition: IF  $\phi: E' \rightarrow \mathbb{R}$ is linear and continuous for the $\sigma(E', E)$ topology then
there exists $x \in E$ such that  $\phi = J_E(x)$.
Set $W = \{f \in E': |\phi(f)| < 1\}$. Since $0_{E'} \in W$ and  $W \in \sigma(E', E)$ there exists $x_1,\dots,x_n \in E$, $\epsilon > 0$ such that if
we set $V_{x_i}(0_{E'},\epsilon) = \{f \in E': |J_E(x_i) - J_E(x_i)(0_{E'})| < \epsilon\}$ then
\[
    \bigcap_{i=1}^{n}V_{x_i}(0_{E'},\epsilon) \subset W\}
\] 
where above the set is
\[
    \{f \in E' : |f(x_i)| < \epsilon\}
\] 
if $f \in \bigcap_{i=1}^{n}Ker(J_E(x_i))$ then for every $\lambda \in R$, $\lambda f \in \bigcap_{i=1}^{n}Ker(J_E(x_i)) \subset \bigcap_{i=1}^{n}V_{x_i}(0_{E'}, \epsilon) \subset W$
And so $|\phi(\lambda f)| < 1$  for all $\lambda \in \mathbb{R}$ and so, $\phi(f) = 0$ in conclusion, we proved that
\[
    \bigcap_{i=1}^{n}Ker(J_E(x_i)) \subset Ker \phi
\] 
and so, there exists  $\lambda_1,\dots,\lambda_n \in \mathbb{R}$ such that $\phi = \sum_{i=1}^{n}\lambda_i J_E(x_i) = J_E(\sum_{i=1}^{n}\lambda_i x_i)$\\
Reminder: $H \subset E'$ is a hyperplane  if there exists $\phi: E' \rightarrow \mathbb{R}$ linear and $\tau_{E'}$ -continuous,  $\alpha \in \mathbb{R}$
such that $\phi \ne 0_{E''}$  and
\[
    H = \{ f \in E' : \phi(f) = \alpha\}
\] 
note that $H^{c} \in \tau_{E'}$
 
\section{Theorem}
let $\phi : E' \rightarrow \mathbb{R}$ be a linear map such that $\phi \ne 0$ identically.
if $\alpha \in \mathbb{R}$ is such that
\[
  H :=  \{f \in E' : \phi(f) = \alpha\}
\] 
is closed for the $\sigma(E', E)$ topology then there exists $x \in E$ such that $\phi = J_E(x)$  $\phi(f) = f(x)$.\\
Proof: Set $W = H^{c}$ then $W \ne \emptyset$. and so we can choose  $f_0 \in W$ since $W \in \sigma(E', E)$
there exists $x_1,\dots,x_n \in E$ and $\epsilon > 0$ such that
\[
   V :=  \bigcap_{i=1}^{n}\{f \in E' : |f(x_i) - f_0(x_i)| < \epsilon\} =  \bigcap_{i=1}^{n}\{f \in E' : |J_E(x_i)(f) - J_E(x_i)(f_0)| < \epsilon\} \subset W = \{f \in E' : \phi(f) \ne \alpha\}
\] 
Note that $V$ is a convex set and so, if $f_1,f_2 \in V$ $\{(1-t)f_1 + tf_2 : t \in [0,1]\} \subset V$. Since $t \rightarrow \phi((1-t)f_1  + tf_2)$ is continuous,
either $\phi(f_1), \phi(f_2) < \alpha$ or  $\phi(f_1),\phi(f_2) > \alpha$. This measn either $V \subset \{\phi < \alpha\}$ or $V \subset \{\phi > \alpha\}$.
Assume for instance that $V \subset \{\phi < \alpha\}$. Note that $U := V - f_0 \in \sigma(E', E)$
and 
\[
    U = \bigcap_{i=1}^{n}\{g \in E' : |g(x_i| < \epsilon\}
\] 
But $U$ can pass over $\alpha$?
If $g \in U$ then  $-g \in U$ equivalently  $g + f_0 \in V \subset W$ and so $\phi(g + f_0) < \alpha$ and $-g + f_0 \in V$ and so $\phi(-g + f_0) < \alpha$
\[
\phi(g) < \alpha - \phi(f_0) \quad -\phi(g) < \alpha - \phi(f_0) \implies |\phi(g)| < \alpha - \phi(f_0)
\] 
so
\[
|\phi(f-f_0)| < \alpha - \phi(f_0) \qquad f \in V
\] 
hence $\phi$ is bounded on $V$.
Exercise: Show that if there exists a neighborhood $O \in \sigma(E', E)$ such that $0_{E'} \in O$ and  $\phi(O)$ is bounded
then $\phi$ is continuous at $0_E'$

\section{Wednesday}
Let $E$ be a Banach space and recall that $J_E : E \rightarrow E''$ is given by $J_E(x)(f) = f(x)$ for all  $x \in E$ and  $f \in E'$.
And Reminder:
 \begin{enumerate}
     \item Let $\alpha \in \mathbb{R}, \phi \in E''$, $\phi \ne 0$ and let $H_{\alpha} = \{\phi = \alpha\}$. If  $H_\alpha$ is closed
         for the $\sigma(E', E)$ topology then $\phi \in \text{range}(J_E)$
     \item Let $C \subset E'$ be a nonempty convex set, then $C$ is closed for  $\tau_{E'}$ iff,  $C$ is closed for  $\sigma(E',E'')$
\end{enumerate}
Summary Assume that $E$ is not reflexive: $E'' \ne J_E(E)$. Let  $\phi \in E'' \setminus J_E(E)$.
Then $Ker(\phi)$ is closed the $\tau_{E'}$-topology an d so it is closed for the  $\sigma(E', E'')$ topology.
By $(a)$, Ker( $\phi$) is not closed for the $\sigma(E', E)$ topology.

\section{PRODUCT TOPOLOGIES}
Let $I \ne \emptyset$ and suppose that for each  $\alpha \in I$ we have a topological space  $( \mathcal{X}_\alpha, \tau_\alpha)$. We define the product space
$\Pi_{\alpha \in I} \mathcal{X}_\alpha$ by
\[
    \{ \omega : I \rightarrow \bigcup_{\alpha \in I} \mathcal{X}_\alpha: w(\alpha) \in \mathcal{X}_\alpha \forall \alpha \in I\}
\] 
we define the projection $\pi_\beta : \Pi_{\alpha \in I} \mathcal{X}_\alpha \rightarrow \mathcal{X}_\beta $ also represented as $proj_\beta$  by
 \[
\pi_\beta(\omega) = \omega(\beta)
\] 
\emph{
    Notation: If $A_\alpha \subset \mathcal{X}_\alpha$ for $\alpha \in J$ where $J \subset I$, we define
     \[
         \Pi_{\alpha \in J} A_\alpha = \{\omega \in \mathcal{X}: \omega(\alpha) \in A_\alpha \forall \alpha \in J\} = \Pi_{\alpha \in I} B_\alpha
    \] 
    where
    \[
    B_\alpha = 
    \begin{cases}
        A_\alpha \qquad \alpha \in J\\
        \mathcal{X}_\alpha \qquad \alpha \not\in J
    \end{cases}
    \] 
}
\\
\emph{
    Definition: the product topology $\tau_{ \mathcal{X}}$ on $ \mathcal{X}$ is the coarsest topology for which all the projections $\pi_\alpha : \mathcal{X} \rightarrow \mathcal{X}_\alpha$ are continuous.
}
\\
Remark:  A base for $\tau_{ \mathcal{X}}$ is the set
\[
    \Pi_{\alpha \in J} A_\alpha
\] 
where $J \subset I$ is of finite cardinality and $A_{\alpha} \subset \mathcal{X}_\alpha$ is open.\\
In the sequel we will assume that $I = E$ and  $ \mathcal{X}_\alpha = \mathbb{R}$ and the $\Pi_{\alpha \in E} X_\alpha = \mathbb{R}^{E} \supset E'$ 
We endow $E'$ with the  $\sigma(E', E)$ topology.
Define $\Phi_E: E' \to \mathbb{R}^{E}$ by $\Phi_E(f) = f$ for  $f \in  E'$ 

\section{Exercise}
\begin{enumerate}
    \item Show that $\Phi_E$ is continuous
    \item Show that $\Phi_E$ is one-one, and $\Phi_E^{-1}$ is continuous.
\end{enumerate}
Solution: It suffices to show that for every $x \in E$ $\pi_x \circ \Phi_E : E' \rightarrow \mathbb{R}$ is continuous.
But $\pi_X \circ \phi_E(f) = f(x)$. Since $f \rightarrow f(x)$  is $\sigma(E',E)$ continuousa We conclude the proof.
Clearly $\Phi_E$ is one-one. It remains to show that $(\Phi_E)^{-1}$ is continuous, We are to show that If $O \in \sigma(E',E)$ then
$\Phi_E(O)$ is open  in  $Y = \Phi_E(E')$ This amounts to showing that  $\Phi_E(O) = V \cap Y$ where  $V \subset \mathbb{R}^{E}$ is open.
Assume that there exists $x_1,\dots,x_n \in E$, $f_0 \in E'$ and $\epsilon > 0$ such that
\[
    O = \bigcap_{i=1}^{n}\{f \in E': |f(x_i) - f_0(x_i)| < \epsilon\}
\] 
then
\[
    \Phi_E(O) = \bigcap_{i=1}^{n}\{\Phi_E(f) : |\Phi_E(f)(x_i) - f_0(x_i)| < \epsilon\} = \bigcap_{i=1}^{n}\{\Phi_E(f) : |\pi_{x_i}\Phi_E(f) - f_0(x_i)| < \epsilon\}
\] 
\[
    = \bigcap_{i=1}^{n}\{ \omega \in \mathbb{R}^{E} : |\pi_{x_i}(\omega) - f_0(x_i)| < \epsilon\} \cap \Phi_E(E')
\] 
\[
    = \bigcap_{i=1}^{n} (\pi_{x_i})^{-1}(f_0(x_i) - \epsilon, f(x_i) + \epsilon) \cap \Phi_E(E')
\] 

\section{Friday}
General Topology\\
Def: Let $( \mathcal{X}, \tau)$ be a topological space
\begin{enumerate}
    \item We say that $ \mathcal{X}$ is compact if every open cover of $ \mathcal{X}$ has a finite subcover
    \item  We say that $ \mathcal{X}$ is sequentially compact  if every sequence in $ X$ admits a
        converges subsequence
    \item We say that   $ \mathcal{X}$ is totally bounded if $ \mathcal{X}$ is complete, and for every $\epsilon > 0$
        there exists  $x_1,\dots,x_n$ such that $X = \bigcup_{i=1}^{n}D_\epsilon(x_i)$ where $ \mathcal{X}$ is a metric space.
\end{enumerate}
Warning, sequentially compact does not imply compact, compact implies sequentially compact. Only on metric spaces, sequentially compact iff compact.\\
Definition: $ (\mathcal{X}, \tau)$ satisfies the finite intersection property,if the following holds:
whenever $(F_i)_{i \in I}$ is a collection of closed sets such that  $F_{i_1} \cap \dots \cap F_{i_n} \ne \emptyset$  for all $n \in \mathbb{N}$ and $i_1,\dots,i_n \in I$, then
\[
    \bigcap_{i\in I} F_i \ne \emptyset
\] 
Finite intersection property if and only if compactness.
Assume $ \mathcal{X}$ is compact and $(F_i)_{i \in I}$ is a collection of closed sets in  $ \mathcal{X}$ such that $\bigcap_{i \in J} F_i \ne \emptyset $ if  $ J \subset I$ finite cardinality.
Assume on the contrary that $\bigcap_{i \in I} F_i = \emptyset$ then $\bigcup_{i \in I}F_i^{c} = X$ and since $X$ is compact and $F_i^{c}$ open there exists  $J \subset I$ finite cardinality
such that
$\bigcup_{i \in J} F_i^{c} = \mathcal{X} \implies \bigcap_{i \in J}F_i = \emptyset$ a contradiction. The reverse direction is similarly obvious.
so we have shown that $FIP \iff cpt$. Indeed assume that  $(O_i)_{i \in I}$ is open cover then  $\bigcap_{i \in I} O_i^{c} = \emptyset$ which implies that there exists finite set $J \subset I$ such that $\bigcap_{i \in J} O_i^{c} = \emptyset \implies \bigcup_{i \in J} O_i = \mathcal{X}$.
\\
Exercise: Let $ \mathcal{X}$ be a metric space. Then the following are equivalent
\begin{enumerate}
    \item $ \mathcal{X}$ is compact
    \item $ \mathcal{X}$ is totally bounded and complete
\end{enumerate}
Solution: Assume $i$, let  $\epsilon > 0$. Then $ \mathcal{X} = \bigcup_{x \in \mathcal{X}} D_\epsilon(x)$ 
since $ \mathcal{X}$ is compact there exists $x_1,\dots,x_n$ such that $ \mathcal{X} = \bigcup_{i=1}^{n}D_\epsilon(x_i)$ thus $ \mathcal{X}$ is totally bounded.
For completeness, let $ (x_n)_n \subset \mathcal{X}$ be a cauchy sequence. Since $ \mathcal{X}$ is sequentially compact
there exists a subsequence $(x_{n_k})_k$ converging to  $x$ hence $(x_n)_n$ converges to  $ x$.
Assume $ii$, It suffices to show that $ \mathcal{X}$ is sequentially compact. Let $(x_n)_{n \in \mathbb{N}} \subset \mathcal{X}$.
Choose $y_1^{1},\dots,y_{k_1}^{1} in \mathcal{X}$  such that $ \mathcal{X} = \bigcup_{i=1}^{k_1}D_1(y_i^{1})$ then at least one of the balls contains infinitely many $(x_n)_n$ say  $D_1(x_1)$ set $S_1 = \{n \in \mathbb{N}: X_n \in D_1(x_1^{1})\}$ there are infinitely many numbers in this set.
We repeat the operation to find $y_k \in X$ such that
\[
    S_k = \{n : S_{k-1} : x_n \in D_{\frac{1}{k}}(y_k)\}
\] 
set of infinite cardinality inductively.a choose $(x_{n_k})$ such that $n_k \in S_k$ and  $n_{k+1} > n_k$. Note that  $(x_{n_k})_k$ is a cauchy
subsetquence of  $(x_n)_n$ and so  $(x_{n_k})$ converges. You can replace this by bounded and closed condition if $X$ is a subset of a finite dimensional vector space.\\
Assumption : $ E$ is a Banach space,  $\phi_E : E' \rightarrow \mathbb{R}^{E}$, $\phi_E(f)(x) = f(x)$ and $\Phi_E$ is one-one
 \[
     \phi_E^{-1} \text{ is continuous if $ \mathbb{R}^{E} = \Pi_E \mathbb{R}$}
\] 
and $E'$ is endowed with the  $\sigma(E', E)$ topology.\\
Exercise: Set $K = \{\omega \in \mathbb{R}^{E}: |\omega(x)| \le ||x||_E,  = \omega(x) + \lambda\omega(y) \forall x,y \in E, \lambda \in \mathbb{R}\}$
Show that $\Phi_E(B_E'(1)) = K$.
\section{theorem}
Theorem: Banach-Alaoglu-Bourbaki:
The set  $B_E'(1)$ is compact for the  $\sigma(E', E)$ topology.\\
Proof: Since $K = \Phi_E(B_{E'}(1))$ it suffices to show that  $K \subset \mathbb{R}^{E}$ is compact 
for the product topology. Set
\[
    K_1 = \Pi_{x \in E}[-||x||_E, ||x||_E]
\] 
By Tychonoff ( the product of compact sets is compact) $K_1$ is compact. For each $x \in E$,  $w \rightarrow w(x)$ is continuous and so 
\[
w \rightarrow w(x+y) - w(x) - w(y), \qquad w \to w(\lambda x) - \lambda w(x)
\] 
are continuous Thus
\[
    K_2 = \{w \in \mathbb{R}^{E}: w(x+y) - w(x) - w(y) = 0, w(\lambda x) - \lambda(x) = 0\}
\] 
is closed in $ \mathbb{R}^{E}$ Since $K = K_1 \cap K_2$ and $K_1$ is compact and $K_2$ is losed we conclude that $K$ is compact.


\end{document}

